% =============================================================================
% 3.4 Analyse comparative
% =============================================================================
\section{Analyse comparative}
\label{sec:ch3-comparaison}

\subsection{Approches de recherche}

Le tableau~\ref{tab:perf-recherche} se relit comme une confirmation de la
complémentarité argumentée au chapitre~\ref{chap:etat-art}. BM25 seul est
le plus rapide et le plus faible en précision~: il retrouve les
identifiants, mais rate les reformulations. Le dense seul inverse le
profil. La fusion récupère le meilleur des deux classements. Le
ré-ordonnancement n'améliore que faiblement le rappel --- les extraits
pertinents étaient déjà dans les cinquante candidats --- mais améliore
nettement la tête de liste, ce qui est précisément ce que le générateur
consomme.

\subsection{Approches de découpage}

\begin{btptable}{Comparaison des approches de découpage}
\label{tab:comp-chunking}
\begin{tabularx}{\linewidth}{@{}L{3.8cm} C{3.0cm} C{3.0cm}@{}}
\toprule
\rowcolor{btppale}
\btpth{Approche} & \btpth{Ancrage} & \btpth{Exhaustivité} \\
\midrule
Fenêtre fixe          & $0{,}65$ & $0{,}55$ \\
Découpage structurel  & $0{,}78$ & $0{,}70$ \\
Parent-enfant         & $0{,}82$ & $0{,}78$ \\
\bottomrule
\end{tabularx}
\end{btptable}

L'hypothèse H2 est soutenue. Le passage de la fenêtre fixe au
parent-enfant gagne $0{,}17$ d'ancrage, soit plus de vingt-cinq pour cent
relatifs, au-delà de la cible de quinze pour cent. Le découpage
structurel, sans hiérarchie parent-enfant, capte déjà une grande part du
gain~: respecter les articles vaut mieux que découper au hasard. Le
parent-enfant ajoute le reste, en donnant au modèle un contexte plus
large que l'unité de recherche.

\subsection{Approches de génération}

\begin{btptable}{Comparaison des approches de génération}
\label{tab:comp-generation}
\begin{tabularx}{\linewidth}{@{}L{4.2cm} C{2.6cm} C{2.4cm} C{2.2cm}@{}}
\toprule
\rowcolor{btppale}
\btpth{Approche} & \btpth{Conformité} & \btpth{Qualité} & \btpth{Versions} \\
\midrule
Tout en un                 & $0{,}55$ & $0{,}60$ & Non \\
Section par section        & $0{,}78$ & $0{,}75$ & Oui \\
Section et revue humaine   & $0{,}92$ & $0{,}88$ & Oui \\
\bottomrule
\end{tabularx}
\end{btptable}

L'hypothèse H4 est nettement soutenue~: la génération par section gagne
$0{,}23$ de conformité sur la génération globale. L'hypothèse H5 l'est
aussi~: la revue humaine porte la conformité à $0{,}92$, au prix d'un
allongement du temps de cycle. Ce surcoût est acceptable --- et, dans le
domaine, non négociable --- dès lors que l'offre engage l'entreprise.

\begin{btpbox}[Lecture croisée]
Les trois comparaisons racontent la même histoire. Chaque fois que l'on
remplace une opération globale et opaque par une opération locale,
contrôlée et, le cas échéant, interrompue, la qualité mesurable
s'améliore. L'architecture en entonnoir, le parent-enfant et la revue
humaine sont trois instances de ce même principe.
\end{btpbox}
